% --- Encoding and Fonts ---
\usepackage[utf8]{inputenc}
\usepackage[T1]{fontenc}
\usepackage{microtype}  % Improves typography

% --- Geometry ---
\usepackage{geometry}
\geometry{margin=1in}

% --- Core Math Packages ---
\usepackage{amsmath, amssymb, amsthm, amsfonts}
\usepackage{mathtools}   % Extends amsmath
\usepackage{bm}          % Bold math symbols
\usepackage{mathrsfs}    % Script fonts (\mathscr)
\usepackage{dsfont}      % Indicator functions (\mathds{1})

% --- Differential Operators (Alternative to physics) ---
\usepackage{derivative}  % For \odv, \pdv, \dv

% --- Figures and Tables ---
\usepackage{graphicx}
\usepackage{caption}
\usepackage{subcaption}
\usepackage{float}
\usepackage{booktabs}    
\usepackage{array}       % Better column formatting
\usepackage{tabularx}    % 
\usepackage{xcolor}

% --- Lists ---
\usepackage{enumitem}

% --- Units ---
\usepackage{siunitx}

% --- Diagrams ---
\usepackage{tikz}
\usetikzlibrary{arrows, positioning, shapes, decorations.pathreplacing}

% --- Headers and Footers ---
\usepackage{fancyhdr}
\pagestyle{fancy}
\fancyhf{}
\fancyhead[L]{\leftmark}
\fancyhead[R]{\thepage}
\renewcommand{\headrulewidth}{0.4pt}

% --- Spacing ---
\usepackage{setspace}
\onehalfspacing  % or \doublespacing for thesis

% --- Hyperlinks ---
\usepackage{hyperref}
\hypersetup{
    colorlinks=true,
    linkcolor=blue,
    citecolor=blue,
    urlcolor=blue,
    filecolor=blue
}
\usepackage{cleveref}  % Must be loaded after hyperref

% --- Enhanced Theorems ---
\usepackage{thmtools}  % Better control over theorem environments

\declaretheoremstyle[
    headfont=\bfseries,
    notefont=\bfseries,
    bodyfont=\itshape,
    spaceabove=6pt,
    spacebelow=6pt
]{definition}
\declaretheoremstyle[
    headfont=\bfseries,
    notefont=\bfseries,
    bodyfont=\normalfont,
    spaceabove=6pt,
    spacebelow=6pt
]{plain}
\declaretheoremstyle[
    headfont=\bfseries,
    notefont=\bfseries,
    bodyfont=\normalfont,
    spaceabove=6pt,
    spacebelow=6pt
]{remark}

\declaretheorem[style=definition, name=Definition, numberwithin=section]{definition}
\declaretheorem[style=plain, name=Proposition, numberwithin=section]{proposition}
\declaretheorem[style=plain, name=Theorem, numberwithin=section]{theorem}
\declaretheorem[style=definition, name=Example, numberwithin=section]{example}
\declaretheorem[style=remark, name=Remark, unnumbered]{remark}

% --- Custom Commands ---
\newcommand{\R}{\mathbb{R}}
\newcommand{\N}{\mathbb{N}}
\newcommand{\Z}{\mathbb{Z}}
\newcommand{\Q}{\mathbb{Q}}
\newcommand{\C}{\mathbb{C}}
\newcommand{\Prob}{\mathbb{P}}
\newcommand{\E}{\mathbb{E}}
\newcommand{\Leb}{\mathcal{L}}
\newcommand{\Borel}{\mathcal{B}}
\newcommand{\fiber}{F}
\newcommand{\confocal}{\Omega}

% --- Math Operators ---
\DeclareMathOperator{\supp}{supp}
\DeclareMathOperator{\Var}{Var}
\DeclareMathOperator{\Cov}{Cov}
\DeclareMathOperator{\diam}{diam}
\DeclareMathOperator{\dist}{dist}
\DeclareMathOperator{\grad}{grad}
\DeclareMathOperator{\curl}{curl}
\DeclareMathOperator{\div}{div}
\DeclareMathOperator{\tr}{tr}  % trace

% --- Title and Author ---
\title{\textbf{How are We Measuring Difference between Two Functions}}
\author{
    Sarbajit Manna
}
\date{\today}